\pagestyle{fancy}

\section{Conclusion}
\label{sec:conclusions}

This chapter returns to the motivation of Chapter 2 and the research question of Chapter 3, states the conclusion that the evidence of Chapter 6 supports, and develops the implications of that conclusion for alignment practice, interpretability methodology, and future work. The structure mirrors the funnel of the thesis in reverse: Section 8.1 states the conclusion at the level of confidence the data warrants; Section 8.2 returns to the motivation and maps the findings onto the questions that originally framed the project; Section 8.3 draws out the implications for alignment, arguing that a wrapper-consistent picture of instruction tuning changes what an alignment team should expect to find when auditing a base model; Section 8.4 reframes the pipeline itself as a reusable methodological contribution that can be applied beyond this specific model pair; Section 8.5 identifies the single most important piece of future work---behaviourally grounded causal validation---and the shorter list of complementary extensions that follow from the limitations of Chapter 7; and Section 8.6 closes with a brief reflection on what this thesis does and does not claim to have established.

\subsection{The Conclusion, Stated at the Level the Data Supports}
\label{subsec:conclusion-stated-at-supported-level}

The evidence assembled in Chapter 6 supports the following conclusion, stated in descriptive and correlational terms. At the residual stream of layer 13 in Gemma 3 1B, the features most behaviourally relevant to instruction-following and safety-related prompts are predominantly reorganised pre-existing features rather than newly introduced ones. The strongest version of this statement---that instruction tuning introduces no new features at all---is not supported by the evidence: the IT-exclusive class identified by the classification of \S6.1 is not empty, and a non-trivial fraction of active IT features lack a strong counterpart in the PT dictionary under the primary operating point thresholds. What the evidence does support is the weaker and more defensible claim that the features that carry the behavioural signatures of instruction tuning---responding to instruction-shaped prompts and to harmful content---are not found in the IT-exclusive class but in the modified class, whose members by construction have counterparts in the base model. The IT-exclusive class exists; it is simply not where the behavioural action is located.

This conclusion is consistent with the wrapper hypothesis of \textcite{jain_mechanistically_2024}, which was originally formulated in synthetic procedurally defined task settings using pruning and linear probing as the interpretability tools. The present work extends that hypothesis, for the first time to the author's knowledge, to a real instruction-tuned language model at the feature level, using Sparse Autoencoders as the descriptive vocabulary and using targeted differential analysis on three behavioural probe datasets as the behavioural grounding. The extension is not a direct confirmation in the causal sense---Chapter 7 is explicit about the absence of a behaviourally grounded causal validation and about the limited reach of the cross-SAE robustness check reported in \S7.3---but it is a substantive descriptive result that links the Jain et al.\ proposal to a production-class model and a production-relevant task setting. Two findings bear the weight of this conclusion and are worth re-identifying in the conclusion by reference rather than by restatement: the cross-model feature-index overlap observed in the instruction-following contrast of \S6.3, where top-ranked differential features in the IT model coincide with top-ranked differential features in the PT model and have order-of-magnitude effect sizes in both; and the distributional pattern of \S6.2 and \S6.1, where instruction-style features are concentrated in the modified class rather than either the preserved or the IT-exclusive class.

The conclusion is also consistent with the broader line of mechanistic fine-tuning work that has accumulated since Jain et al.: \textcite{prakash_fine-tuning_2024} report that fine-tuning enhances existing mechanisms in a case study on entity tracking rather than installing new ones, and the sparse-model-diffing literature that has emerged more recently has converged on similar findings in other model families. This thesis contributes a new data point to that emerging picture, specifically at the feature level, in a real instruction-tuned LLM, and using matched-SAE comparison as the primary tool. The agreement across independent settings and independent tools is in itself weak but non-trivial evidence that the wrapper reading of fine-tuning is not an artefact of any single methodology.

\subsection{Return to the Motivation and the Research Question}
\label{subsec:return-to-motivation-and-rq}

Chapter 2 motivated this thesis on the grounds that post-training is where most user-facing language model behaviour comes from, and that the mechanism of instruction tuning is contested between two broad readings: a ``superficial alignment'' reading in which fine-tuning primarily adjusts style and format over a substrate of base-model capabilities \parencite{zhou_lima_2023}, and a ``new capability'' reading in which fine-tuning can install genuinely new computational structure. Chapter 3 narrowed this motivation to a specific, answerable research question: at the feature level, does instruction tuning introduce new features into a language model's internal representations, or does it reorganise existing ones? The pipeline of Chapter 5 and the results of Chapter 6 were designed to answer exactly this question, and the answer the evidence supports---stated in \S8.1---is on the reorganisation side of the dichotomy without going all the way to ``no new features at all.''

This answer connects back to the superficial-alignment literature in a specific and measurable way. Zhou et al.'s original LIMA observation was that a surprisingly small amount of high-quality supervised fine-tuning data can produce strong instruction-following behaviour, and their interpretation was that the substrate for this behaviour must already exist in the base model. The feature-level evidence of this thesis is consistent with that interpretation in a way that earlier work could not be: the substrate is not only implicit in the LIMA-style behavioural observation but is directly visible as specific SAE feature indices that are already selective for instruction-shaped prompts in the base model and whose selectivity is amplified by instruction tuning. The feature indices identified in \S6.3---feature 2006, feature 446, and the others in the cross-model overlap set---are concrete representational objects that the superficial-alignment literature argues must exist, shown here to actually exist and to be individually identifiable.

The safety contrast tells a weaker but parallel version of the same story, and it matters specifically for the alignment community's concern with how safety behaviours arise through post-training. The finding of \S6.3 that the PT model's top differential features on harmful prompts overlap substantially with the IT model's, and that no feature in the IT top-20 safety list belongs to the IT-exclusive class, implies that the representational basis for the safety-relevant differential signal is not installed by instruction tuning but is already present in the base model as selectivity of knowledge-semantics and instruction-style features for harmful content. Instruction tuning reweights which of these base-model features are most active on harmful prompts, but it does not appear to create a dedicated safety-features cluster from nothing. This is a narrower and more specific version of the claim in Jain et al.\ that fine-tuning installs a ``minimal transformation \ldots\ on top of the underlying model capabilities'' \parencite{jain_mechanistically_2024}: the minimal transformation here is a reweighting of activation patterns over a pre-existing feature basis rather than a change to the basis itself.

\subsection{Implications for Alignment Practice}
\label{subsec:implications-for-alignment-practice}

If the descriptive picture of \S8.1 holds more broadly than this single case study---which is a hypothesis this thesis motivates but cannot prove---then several implications for alignment practice follow directly. They are stated here as hypotheses about what a wrapper-consistent world would imply, not as claims established by the current evidence alone.

\textbf{Base-model auditing becomes a meaningful safety intervention.} In a world where instruction tuning primarily reorganises existing features rather than installing new ones, the base model's feature inventory sets an upper bound on what the instruction-tuned model can represent. A feature that does not exist in some form in the base model cannot be reweighted into existence by instruction tuning, and conversely, a feature that does exist in the base model---whether dormant or merely latent under the base-model training distribution---becomes a candidate for amplification under fine-tuning. The practical consequence is that an alignment team concerned about what behaviours an instruction-tuned model might exhibit has a principled reason to audit the base model's SAE features for latent capabilities that fine-tuning might later surface. This is a specific version of the concern raised in the jailbreak literature \parencite{zou_universal_2023} that behaviours absent from the instruction-tuned model's typical outputs may still be recoverable through adversarial prompting, and a specific version of the concern raised in the sleeper-agent literature \parencite{hubinger_sleeper_2024} that safety training may mask rather than remove harmful capabilities. The feature-level evidence of this thesis supports both concerns in a narrower form: the capabilities that fine-tuning operates on are already in the base model's feature inventory, so any adversarial prompting that bypasses fine-tuning should be expected to surface base-model features, not to invoke instruction-tuning-specific ones.

\textbf{Safety fine-tuning is a reweighting, not a replacement.} The \S6.3 finding that the IT model's top differential safety features are substantially drawn from the same base-model feature population implies that safety behaviours in instruction-tuned models are carried by modulations of pre-existing features rather than by a dedicated safety circuit added during fine-tuning. This is consistent with the observation in the safety-fine-tuning literature \parencite{jain2024makesbreakssafetyfinetuning} that safety fine-tuning can be undone by further fine-tuning on superficially unrelated tasks: if safety is a reweighting of existing features rather than an additive structural change, then any subsequent training that perturbs those weights---in directions that need not be adversarial or safety-targeted---can plausibly disrupt the reweighting. The feature-level picture here does not prove this mechanism, but it is the kind of feature-level evidence that such a mechanism would imply.

\textbf{Cross-checkpoint feature transfer is more principled than arbitrary-direction transfer.} The small robustness check reported in \S7.3 found that SAE decoder directions learned on one checkpoint produce structurally distinguishable effects from random directions of matched magnitude when injected into the other checkpoint. This is a weak finding in the sense that it does not validate any specific behavioural steering claim, but it is a non-trivial finding in the sense that it suggests the feature basis learned by a Gemma Scope 2 SAE is not arbitrary from the perspective of the model---it corresponds to directions the model processes differently from random perturbations. For practitioners who might consider using SAE features trained on one checkpoint to steer, probe, or monitor another closely related checkpoint, this structural observation is weak permission to do so, qualified by the strong caveat that actual behavioural steering claims require their own, task-specific validation.

\subsection{The Pipeline as a Reusable Methodological Contribution}
\label{subsec:pipeline-as-methodological-contribution}

Beyond its specific empirical findings, this thesis contributes a pipeline that can in principle be applied to any model pair for which matched Sparse Autoencoders are available. The pipeline has four stages---activation collection, four-way classification with threshold sensitivity, dual-LLM consensus labelling, and targeted differential analysis on behavioural probes---and each stage is designed to be reused independently rather than only as a component of the whole. A researcher interested in only the classification would need only the activation collection and the Phase 2 matching code; a researcher interested in only the behavioural analysis could apply Phase 4 to any set of features identified by any other means; and the dual-LLM consensus labelling of Phase 3 is model-agnostic and can be applied to any set of SAE features with top-activating contexts.

Two specific methodological choices are worth highlighting as contributions in their own right. First, the four-way classification into preserved, modified, and variant-exclusive classes using two complementary similarity measures (decoder cosine and activation correlation) is more informative than classifications based on a single measure, and the specific observation that the modified class dominates across the entire threshold sensitivity grid is the classification's main empirical yield---it tells a reader that the dominant mode of change between checkpoints is neither replacement nor preservation but gradual reshaping of an existing feature basis. This is the kind of distributional framing that a single-threshold single-measure classification cannot produce. Second, the dual-LLM labelling with structural agreement analysis treats annotator disagreement as a first-class observable rather than as noise to be resolved away, and the specific finding that disagreement concentrates at identifiable taxonomic boundaries (most prominently \texttt{language\_syntax} $\leftrightarrow$ \texttt{instruction\_style}) turns an apparent weakness of automated labelling---that different annotators label the same feature differently---into a measurement of how fuzzy the taxonomic boundaries actually are. Any future SAE-labelling work that runs a single annotator would miss this observation entirely.

These two methodological choices together define a template for matched-SAE comparison studies that future work can apply to other model pairs as those pairs become available. The Gemma Scope 2 release made the Gemma 3 1B comparison possible; a future release that provides matched SAEs for a larger model, for an RLHF-trained rather than SFT-trained variant, or for a non-English model would immediately enable the same pipeline to be re-run and would produce a data point directly comparable to the present results. The main methodological contribution of this thesis, read this way, is not the specific Gemma 3 1B finding but the template that made the finding possible and that can be re-used as matched SAEs become more widely available in the interpretability community.

\subsection{Future Work}
\label{subsec:future-work-directions}

Five directions of future work are motivated directly by the findings and limitations of this thesis, listed in roughly decreasing order of importance.

\textbf{Behaviourally grounded causal validation.} The most important limitation of the current work, as stated at the end of Chapter 7, is the absence of a behaviourally grounded causal experiment that would complement the descriptive evidence with a direct test of whether candidate features actually drive the behaviours they are selective for. A full version of this experiment would steer individual candidate features---for example, the cross-model-shared instruction-style features 2006 and 446, or the cross-model-shared safety-relevant features 262 and 2184---and would measure a task-specific behavioural readout under steering: refusal rate on JailbreakBench prompts under candidate-safety-feature steering, or instruction-following accuracy on Alpaca prompts under candidate-instruction-feature steering. The KL-based readout used in the robustness check of \S7.3 is not adequate for this purpose precisely because it measures distributional shift without measuring whether the shift is in a task-meaningful direction. This experiment is the single most valuable extension of the current pipeline and should be prioritised above all other future work.

\textbf{Multi-layer replication.} Chapter 6 reports results only at layer 13. The classification, labelling, and targeted analysis can in principle be re-run at every layer of the Gemma 3 1B model for which Gemma Scope 2 releases a matched SAE pair, and doing so would produce a layer-by-layer picture of where instruction tuning exerts its effects. The layer-profile of preservation, modification, and exclusive counts is itself a substantive question, and the answer could distinguish between hypotheses in which instruction tuning acts roughly uniformly across the residual stream, hypotheses in which it concentrates at middle layers (where most semantic features are), and hypotheses in which it concentrates at output layers (where most behavioural decisions are made).

\textbf{Cross-family replication.} The findings of this thesis are a single case study on a single model pair. Replicating the pipeline on other model families as matched SAEs become available would test whether the wrapper-consistent picture holds beyond the specific training choices of Google DeepMind's instruction tuning of Gemma 3 1B, or whether it is specific to this recipe. Llama 3, Mistral, and Qwen are all plausible targets; the bottleneck is matched-SAE availability rather than pipeline complexity.

\textbf{Cross-checkpoint crosscoders.} A specific methodological extension of the present work is to replace the two separate SAEs of \S5.1 with a single crosscoder trained jointly on the residual streams of both the PT and IT models, as in the recent sparse-model-diffing literature. A crosscoder would in principle provide a single feature dictionary shared between the two checkpoints and would sidestep the matched-SAE comparison problem by design. The reason crosscoders were not used in this thesis is that the Gemma Scope 2 crosscoder releases are cross-layer rather than cross-checkpoint, and no suitable PT$\leftrightarrow$IT crosscoder was publicly available at the time of this work. A future release of matched PT$\leftrightarrow$IT crosscoders for Gemma 3 would allow the present classification to be re-derived from a single shared dictionary, and the comparison of separate-SAE and crosscoder results would itself be methodologically interesting.

\textbf{Active-feature threshold sensitivity.} A minor but unfinished piece of the current pipeline is a formal sensitivity analysis over the 0.1\% active-feature cutoff used throughout Chapters 5 and 6. Informal checks during development suggested that the classification proportions are not strongly sensitive to this cutoff within the 0.05\%--0.5\% range, but no formal grid was computed. Extending the sensitivity grid of \S5.5 to include the active-feature threshold as a third axis would close this loose end.

\subsection{Closing Reflection}
\label{subsec:closing-reflection}

This thesis set out to test a contested hypothesis about the mechanism of instruction tuning using a specific, answerable feature-level pipeline, and it reports a descriptive answer consistent with the wrapper reading of fine-tuning: at the residual stream of layer 13 in Gemma 3 1B, the features most behaviourally relevant to instruction-following and safety prompts are predominantly reorganised pre-existing features rather than newly introduced ones. The answer is not causal, it is not multi-layer, it is not multi-model, and it is not the last word on the question. What it is, instead, is a concrete and reproducible data point on the feature-level side of a debate that has so far been conducted mostly at the behavioural and circuit level, and a pipeline that can be re-run as matched Sparse Autoencoders become available for other model pairs. The wrapper hypothesis was originally formulated for synthetic tasks and tested with pruning and probing; this thesis extends it to a production-class instruction-tuned language model and tests it with dictionary learning. The convergence of these independent methods on compatible readings is, in the author's view, the most significant thing this thesis contributes to the literature, and the next steps---causal validation, multi-layer replication, cross-family replication---are the natural continuation of the line of work that this thesis joins.
